\subsection{Homological definition of depth}

Let A be a ring, J an ideal of A, and M an A-module.

DEFINITION 1. The depth of M with respect to J is called and denoted \(\mathrm{prof}_{\mathrm{A}}(\mathrm{J};\mathrm{M})\) , or \(\mathrm{prof}(\mathrm{J};\mathrm{M})\) , the lower bound in \(\mathbf{N}\cup\{+\infty\}\) of the set of integers \(n\) such that \(\mathrm{Ext}_{\mathrm{A}}^{n}(\mathrm{A}/\mathrm{J},\mathrm{M})\) is nonzero.

When the ring A is local, one simply calls the depth of M and denotes \(\text{prof}_{\mathrm{A}}(\mathrm{M})\) or \(\text{prof}(\mathrm{M})\) the depth of M with respect to the maximal ideal \(m_{A}\) of A; one calls the depth of the local ring A the depth of the A-module A.

\(\operatorname{Si} \operatorname{prof}_{\mathrm{A}}(\mathrm{J}; \mathrm{M}) = +\infty\) , one has \(\operatorname{Ext}_{\mathrm{A}}^{i}(\mathrm{A}/\mathrm{J}, \mathrm{M}) = 0\) for all \(i\) . \(\operatorname{Si} \operatorname{prof}_{\mathrm{A}}(\mathrm{J}; \mathrm{M}) = r < +\infty\) , one has \(\operatorname{Ext}_{\mathrm{A}}^{i}(\mathrm{A}/\mathrm{J}, \mathrm{M}) = 0\) for \(i < r\) and \(\operatorname{Ext}_{\mathrm{A}}^{r}(\mathrm{A}/\mathrm{J}, \mathrm{M}) \neq 0\) .

Remarks.--- 1) Suppose that the A-module M is finitely generated and that one has M = JM, that is, Supp(M) ∩ V(J) = ∅ (II, § 4, no. 4, cor. of prop. 18). In this case, prof\_A(J;M) is equal to +∞: indeed, the ideal Ann(M) + J is then equal to A (loc. cit.) and is contained in the annihilator of Ext\_A\^{}i(A/J;M) for all i. One will see below (no. 3, cor. 1 of th. 1) that conversely if the ideal J is finitely generated, prof\_A(J;M) = +∞ implies M = JM.

\begin{enumerate}
\def\labelenumi{\arabic{enumi})}
\setcounter{enumi}{1}
\tightlist
\item
  For \(\mathrm{prof}_{\mathbf{A}}(\mathbf{J};\mathbf{M})\) to be zero, it is necessary and sufficient that \(\mathrm{Hom}_{\mathbf{A}}(\mathbf{A} / \mathbf{J},\mathbf{M})\) be nonzero, that is, that M possess a nonzero element annihilated by J; this is in particular the case when one has \(\mathrm{Ass(M)}\cap \mathrm{V(J)}\neq \emptyset\) . If the ring A is noetherian and the A-module M is finitely generated, the following conditions are equivalent (IV, § 1, no. 4, prop. 8):
\end{enumerate}

\begin{enumerate}
\def\labelenumi{\alph{enumi})}
\item
  \(\mathrm{prof}_{\mathrm{A}}(\mathrm{J};\mathrm{M}) = 0\)
\item
  for every \(x \in J\) , the homothety \(x_{\mathbb{M}}\) is not injective;
\end{enumerate}

\[
\mathrm{c)} \text {   we have   } \operatorname{Ass} (\mathrm{M}) \cap \mathrm{V} (\mathrm{J}) \neq \varnothing .
\]

\begin{enumerate}
\def\labelenumi{\arabic{enumi})}
\setcounter{enumi}{2}
\item
  According to remark 2, for a local ring \(\Lambda\) to be of depth zero, it is necessary and sufficient that there exist a nonzero element \(x\) of A such that \(\mathfrak{m}_{\mathrm{A}}x = 0\) . If A is not a field, an element \(x\) such that \(\mathfrak{m}_{\mathrm{A}}x = 0\) is not invertible, hence belongs to \(\mathfrak{m}_{\mathrm{A}}\) and consequently satisfies \(x^{2} = 0\) . Thus a reduced local ring of dimension \(\geqslant 1\) is of depth \(\geqslant 1\) .
\item
  Let \((\mathbf{M}_{\iota})_{\iota \in I}\) be a family of A-modules. According to A, X, p. 89, prop. 7, one has \(\operatorname{prof}(\mathbf{J};\prod_{\iota \in I}\mathbf{M}_{\iota}) = \inf_{\iota \in I}\operatorname{prof}(\mathbf{J};\mathbf{M}_{\iota})\) .
\end{enumerate}

PROPOSITION 1.--- Let A be a ring, J an ideal of A, and 0 → M' → M → M'' → 0 an exact sequence of A-modules. Set

\[
p ^ {\prime} = \operatorname{prof} (\mathrm{J}; \mathrm{M} ^ {\prime}), p = \operatorname{prof} (\mathrm{J}; \mathrm{M}), p ^ {\prime \prime} = \operatorname{prof} (\mathrm{J}; \mathrm{M} ^ {\prime \prime}).
\]

One is then in one of the following three cases, which are mutually exclusive:

\[
p ^ {\prime} = p \leqslant p ^ {\prime \prime}, \quad p = p ^ {\prime \prime} <   p ^ {\prime}, \quad p ^ {\prime \prime} = p ^ {\prime} - 1 <   p.
\]

Consider the exact sequence of extension modules associated with A/J and with the above exact sequence (A, X, p. 92, th. 2). Excluding the case \(p = p' = p'' = +\infty\) ; there then exists in this sequence a first nonzero module, and the following module is also nonzero. This gives the following three possibilities:

\begin{enumerate}
\def\labelenumi{\alph{enumi})}
\item
  The first nonzero module is \(\operatorname{Ext}_{\Lambda}^{p'}(A/J, M')\) . One then has \(p' = p \leqslant p''\) .
\item
  The first nonzero module is \(\operatorname{Ext}_{\Lambda}^{p}(\mathrm{A} / \mathrm{J},\mathrm{M})\) . One then has \(p = p'' < p'\) .
\item
  The first nonzero module is \(\mathrm{Ext}_{\mathrm{A}}^{p''}(\mathrm{A} / \mathrm{J},\mathrm{M}'')\) . One then has \(p'' + 1 = p' \leqslant p\) .
\end{enumerate}

Remark 5.--- Suppose that one has \(p = p'\) and that the injection \(u: M' \to M\) which occurs in the exact sequence of prop. 1 belongs to J Hom \(_A(M', M)\) . One \(a\) then \(p'' = p - 1\) . Indeed, the hypothesis implies that the map Ext \(_A^i(1_{A/J}, u)\) is zero for every integer \(i\) ; this excludes the case \(a\) ) considered above.

PROPOSITION 2.--- Let A be a ring, J an ideal of A, M an A-module, and N an A-module annihilated by a power of J. Then \(\mathrm{Ext}_{\mathrm{A}}^{i}(\mathrm{N},\mathrm{M}) = 0\) for every integer \(i < \mathrm{prof}_{\mathrm{A}}(\mathrm{J};\mathrm{M})\) .

Suppose first that JN = 0 and argue by induction on the integer \(i < \mathrm{prof}_{\mathrm{A}}(\mathrm{J};\mathrm{M})\) . The assertion is evident for \(i < 0\) . Consider N as an (A/J)-module and choose an exact sequence of (A/J)-modules

\[
0 \rightarrow \mathrm{K} \longrightarrow (\mathrm{A} / \mathrm{J}) ^ {(\mathrm{I})} \longrightarrow \mathrm{N} \rightarrow 0.
\]

From this we deduce an exact sequence of extension modules

\[
\operatorname{Ext} _ {\mathrm{A}} ^ {i - 1} (\mathrm{K}, \mathrm{M}) \longrightarrow \operatorname{Ext} _ {\mathrm{A}} ^ {i} (\mathrm{N}, \mathrm{M}) \longrightarrow \operatorname{Ext} _ {\mathrm{A}} ^ {i} ((\mathrm{A} / \mathrm{J}) ^ {(1)}, \mathrm{M}).
\]

The A-module \(\mathrm{Ext}_{\mathrm{A}}^{i-1}(\mathrm{K},\mathrm{M})\) is zero by the induction hypothesis, and the A-module \(\mathrm{Ext}_{\mathrm{A}}^{i}((\mathrm{A}/\mathrm{J})^{(1)},\mathrm{M})\) is isomorphic to \(\mathrm{Ext}_{\mathrm{A}}^{i}(\mathrm{A}/\mathrm{J},\mathrm{M})^{\mathrm{I}}\) (A, X, p.~89, prop. 7), which is zero by definition of the depth. Consequently we have \(\mathrm{Ext}_{\mathrm{A}}^{i}(\mathrm{N},\mathrm{M}) = 0\) .

Let us pass to the general case and argue by induction on the smallest integer m \textgreater{} 0 such that \(J^{m}N = 0\) . We have just treated the case m = 1. Suppose m \textgreater{} 1 and let \(i < \text{prof}_{\text{A}}(\text{J}; \text{M})\) be an integer. Consider the exact sequence

\[
\operatorname{Ext} _ {\Lambda} ^ {i} (\mathrm{N/JN,M}) \longrightarrow \operatorname{Ext} _ {\Lambda} ^ {i} (\mathrm{N,M}) \longrightarrow \operatorname{Ext} _ {\Lambda} ^ {i} (\mathrm{JN,M})
\]

deduced from the exact sequence \(0 \rightarrow JN \rightarrow N \rightarrow N/JN \rightarrow 0\) . The two extreme modules are zero by the induction hypothesis, since N/JN and JN are annihilated by \(J^{m-1}\) . One therefore has \(\operatorname{Ext}_{\Lambda}^{i}(N,M) = 0\) , which is what we wanted to prove.

COROLLARY 1.--- Let m be an integer \textgreater0 and let J' be an ideal of A containing J\^{}m. Then \(\operatorname{prof}_{\mathrm{A}}(\mathrm{J};\mathrm{M}) \leqslant \operatorname{prof}_{\mathrm{A}}(\mathrm{J}';\mathrm{M})\) .

Indeed \(J^{m}\) annihilates the A-module \(A/J'\) , hence \(\operatorname{Ext}_{A}^{i}(A/J', M)\) is zero for every integer \(i < \operatorname{prof}_{A}(J; M)\) (prop. 2).

COROLLARY 2.--- Suppose the ideal J is finitely generated, and let J' be an ideal of A such that V(J) ⊃ V(J').

\begin{enumerate}
\def\labelenumi{\alph{enumi})}
\item
  One \(a \operatorname{prof}_{\mathrm{A}}(\mathrm{J}; \mathrm{M}) \leqslant \operatorname{prof}_{\mathrm{A}}(\mathrm{J}'; \mathrm{M})\) .
\item
  If the ideal J' is finitely generated and if V(J) = V(J'), then we have \(\text{prof}_{A}(J; M) = \text{prof}_{A}(J'; M)\) .
\end{enumerate}

By II, § 4, no. 3, cor. 2 of prop. 11 and § 2, no. 6, prop. 15, there exists an integer m \textgreater{} 0 such that \(J^{m} \subset J'\) . Assertion a) therefore follows from cor. 1, and assertion b) is deduced from it.

Cor. 2 may fail when the ideal J is not finitely generated (exercise 2).

